%-----------------------------------------------------------------------------------------
% Autor dieser Vorlage:
% Stefan Macke (http://fachinformatiker-anwendungsentwicklung.net)
% Permalink zur Vorlage: http://fiae.link/LaTeXVorlageFIAE
%
% Sämtliche verwendeten Abbildungen, Tabellen und Listings stammen von Dirk Grashorn.
%
% Lizenz: Creative Commons 4.0 Namensnennung - Weitergabe unter gleichen Bedingungen
% -----------------------------------------------------------------------------------------

\documentclass[
	ngerman,
	toc=listof, % Abbildungsverzeichnis sowie Tabellenverzeichnis in das Inhaltsverzeichnis aufnehmen
	toc=bibliography, % Literaturverzeichnis in das Inhaltsverzeichnis aufnehmen
	footnotes=multiple, % Trennen von direkt aufeinander folgenden Fußnoten
	parskip=half, % vertikalen Abstand zwischen Absätzen verwenden anstatt horizontale Einrückung von Folgeabsätzen
	numbers=noendperiod % Den letzten Punkt nach einer Nummerierung entfernen (nach DIN 5008)
]{scrartcl}
\pdfminorversion=5 % erlaubt das Einfügen von pdf-Dateien bis Version 1.7, ohne eine Fehlermeldung zu werfen (keine Garantie für fehlerfreies Einbetten!)
\usepackage[utf8]{inputenc} % muss als erstes eingebunden werden, da Meta/Packages ggfs. Sonderzeichen enthalten

% !TEX root = Projektdokumentation.tex

% Hinweis: der Titel muss zum Inhalt des Projekts passen und den zentralen Inhalt des Projekts deutlich herausstellen
\newcommand{\titel}{Parsen eines Schemas in eine Baumstruktur}
\newcommand{\untertitel}{und zergliedern eines Datenstroms anhand dieses Schemas}
\newcommand{\kompletterTitel}{\titel{} -- \untertitel}

\newcommand{\autorName}{René Ederer}
\newcommand{\autorAnschrift}{Steinmetzstr. 2}
\newcommand{\autorOrt}{90431 Nürnberg}

\newcommand{\betriebLogo}{LogoBetrieb.pdf}
\newcommand{\betriebName}{\textsc{Phoenix Group IT GmbH}}
\newcommand{\betriebAnschrift}{Sportplatzstr. 30}
\newcommand{\betriebOrt}{90765 Fürth}

\newcommand{\ausbildungsberuf}{Fachinformatiker für Anwendungsentwicklung}
\newcommand{\betreff}{Dokumentation zur betrieblichen Projektarbeit}
\newcommand{\pruefungstermin}{Sommer 2016}
\newcommand{\abgabeOrt}{Nürnberg}
\newcommand{\abgabeTermin}{15.05.2016}
 % Metadaten zu diesem Dokument (Autor usw.)
% !TEX root = ../Projektdokumentation.tex

% Anpassung an Landessprache ---------------------------------------------------
\usepackage{babel}

% Umlaute ----------------------------------------------------------------------
%   Umlaute/Sonderzeichen wie äüöß direkt im Quelltext verwenden (CodePage).
%   Erlaubt automatische Trennung von Worten mit Umlauten.
% ------------------------------------------------------------------------------
\usepackage[T1]{fontenc}
\usepackage{textcomp} % Euro-Zeichen etc.




% Schrift ----------------------------------------------------------------------
\usepackage{lmodern} % bessere Fonts
\usepackage{relsize} % Schriftgröße relativ festlegen

% Tabellen ---------------------------------------------------------------------
\PassOptionsToPackage{table}{xcolor}
\usepackage{tabularx}
% für lange Tabellen
\usepackage{longtable}
\usepackage{array}
\usepackage{ragged2e}
\usepackage{lscape}
\newcolumntype{w}[1]{>{\raggedleft\hspace{0pt}}p{#1}} % Spaltendefinition rechtsbündig mit definierter Breite

% Grafiken ---------------------------------------------------------------------
\usepackage[dvips,final]{graphicx} % Einbinden von JPG-Grafiken ermöglichen
\usepackage{graphics} % keepaspectratio
\usepackage{floatflt} % zum Umfließen von Bildern
\graphicspath{{Bilder/}} % hier liegen die Bilder des Dokuments

% Sonstiges --------------------------------------------------------------------
\usepackage[titles]{tocloft} % Inhaltsverzeichnis DIN 5008 gerecht einrücken
\usepackage{amsmath,amsfonts} % Befehle aus AMSTeX für mathematische Symbole
\usepackage{enumitem} % anpassbare Enumerates/Itemizes
\usepackage{xspace} % sorgt dafür, dass Leerzeichen hinter parameterlosen Makros nicht als Makroendezeichen interpretiert werden

\usepackage{makeidx} % für Index-Ausgabe mit \printindex
\usepackage[printonlyused]{acronym} % es werden nur benutzte Definitionen aufgelistet

% Einfache Definition der Zeilenabstände und Seitenränder etc.
\usepackage{setspace}
\usepackage{geometry}

% Symbolverzeichnis
\usepackage[intoc]{nomencl}
\let\abbrev\nomenclature
\renewcommand{\nomname}{Abkürzungsverzeichnis}
\setlength{\nomlabelwidth}{.25\hsize}
\renewcommand{\nomlabel}[1]{#1 \dotfill}
\setlength{\nomitemsep}{-\parsep}

\usepackage{varioref} % Elegantere Verweise. „auf der nächsten Seite“
\usepackage{url} % URL verlinken, lange URLs umbrechen etc.

\usepackage{chngcntr} % fortlaufendes Durchnummerieren der Fußnoten
% \usepackage[perpage]{footmisc} % Alternative: Nummerierung der Fußnoten auf jeder Seite neu

\usepackage{ifthen} % bei der Definition eigener Befehle benötigt
\usepackage{todonotes} % definiert u.a. die Befehle \todo und \listoftodos
\usepackage[square]{natbib} % wichtig für korrekte Zitierweise

% PDF-Optionen -----------------------------------------------------------------
\usepackage{pdfpages}
\pdfminorversion=5 % erlaubt das Einfügen von pdf-Dateien bis Version 1.7, ohne eine Fehlermeldung zu werfen (keine Garantie für fehlerfreies Einbetten!)
\usepackage[
    bookmarks,
    bookmarksnumbered,
    bookmarksopen=true,
    bookmarksopenlevel=1,
    colorlinks=true,
% diese Farbdefinitionen zeichnen Links im PDF farblich aus
    linkcolor=AOBlau, % einfache interne Verknüpfungen
    anchorcolor=AOBlau,% Ankertext
    citecolor=AOBlau, % Verweise auf Literaturverzeichniseinträge im Text
    filecolor=AOBlau, % Verknüpfungen, die lokale Dateien öffnen
    menucolor=AOBlau, % Acrobat-Menüpunkte
    urlcolor=AOBlau,
% diese Farbdefinitionen sollten für den Druck verwendet werden (alles schwarz)
    %linkcolor=black, % einfache interne Verknüpfungen
    %anchorcolor=black, % Ankertext
    %citecolor=black, % Verweise auf Literaturverzeichniseinträge im Text
    %filecolor=black, % Verknüpfungen, die lokale Dateien öffnen
    %menucolor=black, % Acrobat-Menüpunkte
    %urlcolor=black,
%
    %backref, % Quellen werden zurück auf ihre Zitate verlinkt
    pdftex,
    plainpages=false, % zur korrekten Erstellung der Bookmarks
    pdfpagelabels=true, % zur korrekten Erstellung der Bookmarks
    hypertexnames=false, % zur korrekten Erstellung der Bookmarks
    linktocpage % Seitenzahlen anstatt Text im Inhaltsverzeichnis verlinken
]{hyperref}
% Befehle, die Umlaute ausgeben, führen zu Fehlern, wenn sie hyperref als Optionen übergeben werden
\hypersetup{
    pdftitle={\titel -- \untertitel},
    pdfauthor={\autorName},
    pdfcreator={\autorName},
    pdfsubject={\titel -- \untertitel},
    pdfkeywords={\titel -- \untertitel},
}


% zum Einbinden von Programmcode -----------------------------------------------
\usepackage{listings}
\usepackage{xcolor}
\definecolor{hellgelb}{rgb}{1,1,0.9}
\definecolor{colKeys}{rgb}{0,0,1}
\definecolor{colIdentifier}{rgb}{0,0,0}
\definecolor{colComments}{rgb}{0,0.5,0}
\definecolor{colString}{rgb}{1,0,0}
\lstset{
    float=hbp,
	basicstyle=\footnotesize,
    identifierstyle=\color{colIdentifier},
    keywordstyle=\color{colKeys},
    stringstyle=\color{colString},
    commentstyle=\color{colComments},
    backgroundcolor=\color{hellgelb},
    columns=flexible,
    tabsize=2,
    frame=single,
    extendedchars=true,
    showspaces=false,
    showstringspaces=false,
    numbers=left,
    numberstyle=\tiny,
    breaklines=true,
    breakautoindent=true,
	captionpos=b,
}
\lstdefinelanguage{cs}{
	sensitive=false,
	morecomment=[l]{//},
	morecomment=[s]{/*}{*/},
	morestring=[b]",
	morekeywords={
		abstract,event,new,struct,as,explicit,null,switch
		base,extern,object,this,bool,false,operator,throw,
		break,finally,out,true,byte,fixed,override,try,
		case,float,params,typeof,catch,for,private,uint,
		char,foreach,protected,ulong,checked,goto,public,unchecked,
		class,if,readonly,unsafe,const,implicit,ref,ushort,
		continue,in,return,using,decimal,int,sbyte,virtual,
		default,interface,sealed,volatile,delegate,internal,short,void,
		do,is,sizeof,while,double,lock,stackalloc,
		else,long,static,enum,namespace,string},
}
\lstdefinelanguage{natural}{
	sensitive=false,
	morecomment=[l]{/*},
	morestring=[b]",
	morestring=[b]',
	alsodigit={-,*},
	morekeywords={
		DEFINE,DATA,LOCAL,END-DEFINE,WRITE,CALLNAT,PARAMETER,USING,
		IF,NOT,END-IF,ON,*ERROR-NR,ERROR,END-ERROR,ESCAPE,ROUTINE,
		PERFORM,SUBROUTINE,END-SUBROUTINE,CONST,END-FOR,END,FOR,RESIZE,
		ARRAY,TO,BY,VALUE,RESET,COMPRESS,INTO,EQ},
}
\lstdefinelanguage{php}{
	sensitive=false,
	morecomment=[l]{/*},
	morestring=[b]",
	morestring=[b]',
	alsodigit={-,*},
	morekeywords={
		abstract,and,array,as,break,case,catch,cfunction,class,clone,const,
		continue,declare,default,do,else,elseif,enddeclare,endfor,endforeach,
		endif,endswitch,endwhile,extends,final,for,foreach,function,global,
		goto,if,implements,interface,instanceof,namespace,new,old_function,or,
		private,protected,public,static,switch,throw,try,use,var,while,xor
		die,echo,empty,exit,eval,include,include_once,isset,list,require,
		require_once,return,print,unset},
}
 % verwendete Packages
% !TEX root = ../Projektdokumentation.tex

% Seitenränder -----------------------------------------------------------------
\setlength{\topskip}{\ht\strutbox} % behebt Warnung von geometry
\geometry{a4paper,left=20mm,right=20mm,top=25mm,bottom=35mm}

\usepackage[
	automark, % Kapitelangaben in Kopfzeile automatisch erstellen
	headsepline, % Trennlinie unter Kopfzeile
	ilines % Trennlinie linksbündig ausrichten]
]{scrpage2}


\usepackage{lastpage}


% Kopf- und Fußzeilen ----------------------------------------------------------
\pagestyle{scrheadings}
% chapterpagestyle gibt es nicht in scrartcl
%\renewcommand{\chapterpagestyle}{scrheadings}
\clearscrheadfoot

% Kopfzeile
\renewcommand{\headfont}{\normalfont} % Schriftform der Kopfzeile
\ihead{\large{\textsc{\titel}}\\ \small{\untertitel} \\[2ex] \textit{\headmark}}
\chead{}
\ohead{\includegraphics[scale=0.2]{\betriebLogo}}
\setlength{\headheight}{15mm} % Höhe der Kopfzeile
%\setheadwidth[0pt]{textwithmarginpar} % Kopfzeile über den Text hinaus verbreitern (falls Logo den Text überdeckt)

% Fußzeile
\ifoot{\autorName}
%\cfoot{}
\cfoot{}
%[\pagemark]{\pagemark{} von \pageref{LastPage}}
\ofoot{Seite \pagemark{} von \pageref{LastPage}}


% Überschriften nach DIN 5008 in einer Fluchtlinie
% ------------------------------------------------------------------------------

% Abstand zwischen Nummerierung und Überschrift definieren
% > Schön wäre hier die dynamische Berechnung des Abstandes in Abhängigkeit
% > der Verschachtelungstiefe des Inhaltsverzeichnisses
\newcommand{\headingSpace}{1.5cm}

% Abschnittsüberschriften im selben Stil wie beim Inhaltsverzeichnis einrücken
\renewcommand*{\othersectionlevelsformat}[3]{
  \makebox[\headingSpace][l]{#3\autodot}
}

% Für die Einrückung wird das Paket tocloft benötigt
%\cftsetindents{chapter}{0.0cm}{\headingSpace}
\cftsetindents{section}{0.0cm}{\headingSpace}
\cftsetindents{subsection}{0.0cm}{\headingSpace}
\cftsetindents{subsubsection}{0.0cm}{\headingSpace}
\cftsetindents{figure}{0.0cm}{\headingSpace}
\cftsetindents{table}{0.0cm}{\headingSpace}


% Allgemeines
% ------------------------------------------------------------------------------

\onehalfspacing % Zeilenabstand 1,5 Zeilen
\frenchspacing % erzeugt ein wenig mehr Platz hinter einem Punkt

% Schusterjungen und Hurenkinder vermeiden
\clubpenalty = 10000
\widowpenalty = 10000
\displaywidowpenalty = 10000

% Quellcode-Ausgabe formatieren
\lstset{numbers=left, numberstyle=\tiny, numbersep=5pt, breaklines=true}
\lstset{emph={square}, emphstyle=\color{red}, emph={[2]root,base}, emphstyle={[2]\color{blue}}}

\counterwithout{footnote}{section} % Fußnoten fortlaufend durchnummerieren
\setcounter{tocdepth}{\subsubsectionlevel} % im Inhaltsverzeichnis werden die Kapitel bis zum Level der subsubsection übernommen
\setcounter{secnumdepth}{\subsubsectionlevel} % Kapitel bis zum Level der subsubsection werden nummeriert

% Aufzählungen anpassen
\renewcommand{\labelenumi}{\arabic{enumi}.}
\renewcommand{\labelenumii}{\arabic{enumi}.\arabic{enumii}.}
\renewcommand{\labelenumiii}{\arabic{enumi}.\arabic{enumii}.\arabic{enumiii}}

% Tabellenfärbung:
\definecolor{heading}{rgb}{0.62,0.76,0.69}
\definecolor{odd}{rgb}{0.87,0.92,0.89}
 % Definitionen zum Aussehen der Seiten
% !TEX root = ../Projektdokumentation.tex

% Abkürzungen, ggfs. mit korrektem Leerraum
\newcommand{\bs}{$\backslash$\xspace}
\newcommand{\bspw}{bspw.\xspace}
\newcommand{\bzw}{bzw.\xspace}
\newcommand{\ca}{ca.\xspace}
\newcommand{\dahe}{\mbox{d.\,h.}\xspace}
\newcommand{\etc}{etc.\xspace}
\newcommand{\eur}[1]{\mbox{#1\,\texteuro}\xspace}
\newcommand{\evtl}{evtl.\xspace}
\newcommand{\ggfs}{ggfs.\xspace}
\newcommand{\Ggfs}{Ggfs.\xspace}
\newcommand{\gqq}[1]{\glqq{}#1\grqq{}}
\newcommand{\inkl}{inkl.\xspace}
\newcommand{\insb}{insb.\xspace}
\newcommand{\ua}{\mbox{u.\,a.}\xspace}
\newcommand{\usw}{usw.\xspace}
\newcommand{\Vgl}{Vgl.\xspace}
\newcommand{\zB}{\mbox{z.\,B.}\xspace}

% Befehle für häufig anfallende Aufgaben
\newcommand{\Abbildung}[1]{\autoref{fig:#1}}
\newcommand{\Anhang}[1]{\appendixname{}~\ref{#1}: \nameref{#1} \vpageref{#1}}
\newcommand{\includegraphicsKeepAspectRatio}[2]{\includegraphics[width=#2\textwidth,height=#2\textheight,keepaspectratio]{#1}}
\newcommand{\Zitat}[2][\empty]{\ifthenelse{\equal{#1}{\empty}}{\citep{#2}}{\citep[#1]{#2}}}
\newcommand{\Autor}[1]{\textsc{#1}} % zum Ausgeben von Autoren
\newcommand{\itemd}[2]{\item{\textbf{#1}}\\{#2}} % erzeugt ein Listenelement mit fetter Überschrift

% fügt Tabellen aus einer TEX-Datei ein
\newcommand{\tabelle}[3] % Parameter: caption, label, file
{\begin{table}[htbp]
\centering
\singlespacing
\input{Tabellen/#3}
\caption{#1}
\label{#2}
\end{table}}

\newcommand{\tabelleAnhang}[1] % Parameter: file
{\begin{center}
\singlespacing
\input{Tabellen/#1}
\end{center}}

% einfaches Wechseln der Schrift, z.B.: \changefont{cmss}{sbc}{n}
\newcommand{\changefont}[3]{\fontfamily{#1} \fontseries{#2} \fontshape{#3} \selectfont}

% Verwendung analog zu \includegraphics
\newlength{\myx} % Variable zum Speichern der Bildbreite
\newlength{\myy} % Variable zum Speichern der Bildhöhe
\newcommand\includegraphicstotab[2][\relax]{%
% Abspeichern der Bildabmessungen
\settowidth{\myx}{\includegraphics[{#1}]{#2}}%
\settoheight{\myy}{\includegraphics[{#1}]{#2}}%
% das eigentliche Einfügen
\parbox[c][1.1\myy][c]{\myx}{%
\includegraphics[{#1}]{#2}}%
}

\definecolor{AOBlau}{rgb}{0, 0.28, 0.56}

% verschiedene Befehle um Wörter semantisch auszuzeichnen ----------------------
\newcommand{\Index}[2][\empty]{\ifthenelse{\equal{#1}{\empty}}{\index{#2}#2}{\index{#1}#2}}
\newcommand{\Fachbegriff}[2][\empty]{\ifthenelse{\equal{#1}{\empty}}{\textit{\Index{#2}}}{\textit{\Index[#1]{#2}}}}
\newcommand{\NeuerBegriff}[2][\empty]{\ifthenelse{\equal{#1}{\empty}}{\textbf{\Index{#2}}}{\textbf{\Index[#1]{#2}}}}

\newcommand{\Ausgabe}[1]{\texttt{#1}}
\newcommand{\Eingabe}[1]{\texttt{#1}}
\newcommand{\Code}[1]{\texttt{#1}}
\newcommand{\Datei}[1]{\texttt{#1}}

\newcommand{\Assembly}[1]{\textsf{#1}}
\newcommand{\Klasse}[1]{\textsf{#1}}
\newcommand{\Methode}[1]{\textsf{#1}}
\newcommand{\Attribut}[1]{\textsf{#1}}

\newcommand{\Datentyp}[1]{\textsf{#1}}
\newcommand{\XMLElement}[1]{\textsf{#1}}
\newcommand{\Webservice}[1]{\textsf{#1}}

\newcommand{\Refactoring}[1]{\Fachbegriff{#1}}
\newcommand{\CodeSmell}[1]{\Fachbegriff{#1}}
\newcommand{\Metrik}[1]{\Fachbegriff{#1}}
\newcommand{\DesignPattern}[1]{\Fachbegriff{#1}}
 % eigene allgemeine Befehle, die z.B. die Arbeit mit LaTeX erleichtern
% !TEX root = ../Projektdokumentation.tex

% Abkürzungen, ggfs. mit korrektem Leerraum
\newcommand{\bs}{$\backslash$\xspace}
\newcommand{\bspw}{bspw.\xspace}
\newcommand{\bzw}{bzw.\xspace}
\newcommand{\ca}{ca.\xspace}
\newcommand{\dahe}{\mbox{d.\,h.}\xspace}
\newcommand{\etc}{etc.\xspace}
\newcommand{\eur}[1]{\mbox{#1\,\texteuro}\xspace}
\newcommand{\evtl}{evtl.\xspace}
\newcommand{\ggfs}{ggfs.\xspace}
\newcommand{\Ggfs}{Ggfs.\xspace}
\newcommand{\gqq}[1]{\glqq{}#1\grqq{}}
\newcommand{\inkl}{inkl.\xspace}
\newcommand{\insb}{insb.\xspace}
\newcommand{\ua}{\mbox{u.\,a.}\xspace}
\newcommand{\usw}{usw.\xspace}
\newcommand{\Vgl}{Vgl.\xspace}
\newcommand{\zB}{\mbox{z.\,B.}\xspace}

% Befehle für häufig anfallende Aufgaben
\newcommand{\Abbildung}[1]{\autoref{fig:#1}}
\newcommand{\Anhang}[1]{\appendixname{}~\ref{#1}: \nameref{#1} \vpageref{#1}}
\newcommand{\includegraphicsKeepAspectRatio}[2]{\includegraphics[width=#2\textwidth,height=#2\textheight,keepaspectratio]{#1}}
\newcommand{\Zitat}[2][\empty]{\ifthenelse{\equal{#1}{\empty}}{\citep{#2}}{\citep[#1]{#2}}}
\newcommand{\Autor}[1]{\textsc{#1}} % zum Ausgeben von Autoren
\newcommand{\itemd}[2]{\item{\textbf{#1}}\\{#2}} % erzeugt ein Listenelement mit fetter Überschrift

% fügt Tabellen aus einer TEX-Datei ein
\newcommand{\tabelle}[3] % Parameter: caption, label, file
{\begin{table}[htbp]
\centering
\singlespacing
\input{Tabellen/#3}
\caption{#1}
\label{#2}
\end{table}}

\newcommand{\tabelleAnhang}[1] % Parameter: file
{\begin{center}
\singlespacing
\input{Tabellen/#1}
\end{center}}

% einfaches Wechseln der Schrift, z.B.: \changefont{cmss}{sbc}{n}
\newcommand{\changefont}[3]{\fontfamily{#1} \fontseries{#2} \fontshape{#3} \selectfont}

% Verwendung analog zu \includegraphics
\newlength{\myx} % Variable zum Speichern der Bildbreite
\newlength{\myy} % Variable zum Speichern der Bildhöhe
\newcommand\includegraphicstotab[2][\relax]{%
% Abspeichern der Bildabmessungen
\settowidth{\myx}{\includegraphics[{#1}]{#2}}%
\settoheight{\myy}{\includegraphics[{#1}]{#2}}%
% das eigentliche Einfügen
\parbox[c][1.1\myy][c]{\myx}{%
\includegraphics[{#1}]{#2}}%
}

\definecolor{AOBlau}{rgb}{0, 0.28, 0.56}

% verschiedene Befehle um Wörter semantisch auszuzeichnen ----------------------
\newcommand{\Index}[2][\empty]{\ifthenelse{\equal{#1}{\empty}}{\index{#2}#2}{\index{#1}#2}}
\newcommand{\Fachbegriff}[2][\empty]{\ifthenelse{\equal{#1}{\empty}}{\textit{\Index{#2}}}{\textit{\Index[#1]{#2}}}}
\newcommand{\NeuerBegriff}[2][\empty]{\ifthenelse{\equal{#1}{\empty}}{\textbf{\Index{#2}}}{\textbf{\Index[#1]{#2}}}}

\newcommand{\Ausgabe}[1]{\texttt{#1}}
\newcommand{\Eingabe}[1]{\texttt{#1}}
\newcommand{\Code}[1]{\texttt{#1}}
\newcommand{\Datei}[1]{\texttt{#1}}

\newcommand{\Assembly}[1]{\textsf{#1}}
\newcommand{\Klasse}[1]{\textsf{#1}}
\newcommand{\Methode}[1]{\textsf{#1}}
\newcommand{\Attribut}[1]{\textsf{#1}}

\newcommand{\Datentyp}[1]{\textsf{#1}}
\newcommand{\XMLElement}[1]{\textsf{#1}}
\newcommand{\Webservice}[1]{\textsf{#1}}

\newcommand{\Refactoring}[1]{\Fachbegriff{#1}}
\newcommand{\CodeSmell}[1]{\Fachbegriff{#1}}
\newcommand{\Metrik}[1]{\Fachbegriff{#1}}
\newcommand{\DesignPattern}[1]{\Fachbegriff{#1}}
 % eigene projektspezifische Befehle, z.B. Abkürzungen usw.

\begin{document}


\phantomsection
\thispagestyle{empty}
\pdfbookmark[1]{Eidesstattliche Erklärung}{ihkdeckblatt}
\includegraphicsKeepAspectRatio{DeckblattIHK}{1}
\cleardoublepage

\phantomsection
\thispagestyle{plain}
\pdfbookmark[1]{Deckblatt}{deckblatt}
% !TEX root = Projektdokumentation.tex
\begin{titlepage}

\begin{center}
\includegraphics[scale=0.25]{LogoIHK.pdf}\\[1ex]
\Large{Abschlussprüfung \pruefungstermin}\\[3ex]

\Large{\ausbildungsberuf}\\
\LARGE{\betreff}\\[4ex]

\huge{\textbf{\titel}}\\[1.5ex]
\Large{\textbf{\untertitel}}\\[4ex]

\normalsize
Abgabetermin: \abgabeOrt, den \abgabeTermin\\[3em]
\textbf{Prüfungsbewerber:}\\
\autorName\\
\autorAnschrift\\
\autorOrt\\[5ex]

\includegraphics[scale=0.3]{\betriebLogo}\\[2ex]
\textbf{Ausbildungsbetrieb:}\\
\betriebName\\
\betriebAnschrift\\
\betriebOrt\\[5em]
\end{center}

\small
\noindent
Dieses Werk einschließlich seiner Teile ist \textbf{urheberrechtlich geschützt}.
Jede Verwertung außerhalb der engen Grenzen des Urheberrechtgesetzes ist ohne
Zustimmung des Autors unzulässig und strafbar. Das gilt insbesondere für
Vervielfältigungen, Übersetzungen, Mikroverfilmungen sowie die Einspeicherung
und Verarbeitung in elektronischen Systemen.

\end{titlepage}
\cleardoublepage
\def \renePage {IV}
% Preface --------------------------------------------------------------------
\phantomsection
\pagenumbering{Roman}
\pdfbookmark[1]{Inhaltsverzeichnis}{inhalt}
\tableofcontents
\cleardoublepage

\phantomsection
\listoffigures
\cleardoublepage

\phantomsection
\listoftables
\cleardoublepage

\phantomsection
\lstlistoflistings
\cleardoublepage

\newcommand{\abkvz}{Abkürzungsverzeichnis}
\renewcommand{\nomname}{\abkvz}
\section*{\abkvz}
\markboth{\abkvz}{\abkvz}
\addcontentsline{toc}{section}{\abkvz}
% !TEX root = Projektdokumentation.tex

% Es werden nur die Abkürzungen aufgelistet, die mit \ac definiert und auch benutzt wurden. 
%
% \acro{VERSIS}{Versicherungsinformationssystem\acroextra{ (Bestandsführungssystem)}}
% Ergibt in der Liste: VERSIS Versicherungsinformationssystem (Bestandsführungssystem)
% Im Text aber: \ac{VERSIS} -> Versicherungsinformationssystem (VERSIS)

% Hinweis: allgemein bekannte Abkürzungen wie z.B. bzw. u.a. müssen nicht ins Abkürzungsverzeichnis aufgenommen werden
% Hinweis: allgemein bekannte IT-Begriffe wie Datenbank oder Programmiersprache müssen nicht erläutert werden,
%          aber ggfs. Fachbegriffe aus der Domäne des Prüflings (z.B. Versicherung)

% Die Option (in den eckigen Klammern) enthält das längste Label oder
% einen Platzhalter der die Breite der linken Spalte bestimmt.
\begin{acronym}[WWWWW]
	\acro{COBOL}{Common Business-Oriented Language}
	\acro{GUI}{Graphical User Interface}
	\acro{IDE}{Integrated Development Environment}
	\acro{SVN}{Subversion}
	\acro{TDD}{Test Driven Development}
	\acro{UML}{Unified Modeling Language}
	\acro{XML}{Extensible Markup Language}
\end{acronym}

\clearpage

% Inhalt ---------------------------------------------------------------------
\pagenumbering{arabic}
\def \renePage {40}
% !TEX root = Projektdokumentation.tex
% !TEX root = ../Projektdokumentation.tex
\section{Projektdefinition}
\label{sec:Einleitung}
\subsection{Auftraggeber}
\label{sec:Auftraggeber}
Die Phoenix Pharmahandel GmbH \& Co. KG. und ihre Tochtergesellschaften sind europaweit unter dem Namen ``Phoenix group'' mit etwa 30000 Mitarbeitern im Pharmagroßhandel tätig. Ihre Haupttätigkeiten sind der Großhandel mit Pharmaprodukten und die Belieferung von Apotheken mit Medikamenten.

Auftraggeber des Projektes ist die Phoenix group IT GmbH. Sie hat etwa 200 Mitarbeiter und unterstützt die Phoenix group durch die Bereitstellung von IT-Dienstleistungen.

\subsection{Projektumfeld} 
\label{sec:Projektumfeld}
Die Phoenix group IT ist weiter unterteilt in die Abteilungen Inbound, Outbound und Warehouse. Das Projekt findet in der Abteilung Warehouse statt, in der überwiegend mit COBOL gearbeitet wird.

Phoenix verwendet ein firmeneigenes Dateiformat als Schnittstelle zu verschiedenen Programmen und Diensten. Es wird im Folgenden 1920Schema genannt und spielt eine zentrale Rolle für das Projekt.

\subsubsection{1920Schemas}
1920Schemas sind Textdateien, die einen Satz hierarchisch gegliederter Variablen beschreiben. Phoenix nutzt diese Schemas unter anderem als Schnittstelle, um Daten vom Mainframe zum Lagerrechner zu schicken, als Vorlage für Copybooks\footnote{COBOL-Datei, in der eine Variablenstruktur definiert wird} und als Schnittstelle zu SSORT\footnote{IBM-Programm, zeigt Copybooks an}.

\textbf{Beispiel für ein 1920Schema}
\tabelleLinks{Beispiel für ein 1920Schema}{tab:Beispiel für ein 1920Schema}{BeispielSchema.tex}

\subsubsection{Datenströme}
Anhand von 1920Schemas werden Datenströme zergliedert und erhalten dadurch eine Bedeutung. Datenströme können prinzipiell beliebige Zeichen\footnote{Ein Byte entspricht einem Zeichen. Die Angaben in 1920Schemas haben die Einheit Byte, im Zusammenhang mit dem Datenstrom ist der Begriff Zeichen aber manchmal anschaulicher. Eine Unterscheidung der beiden Begriffe ist für das Projekt unwesentlich.} enthalten, nur die Anzahl muss ausreichen, um jeder Schema-Variablen einen Wert zuzuweisen. Das Schema aus Tabelle 1 beschreibt einen 17 Byte langen Datenstrom (5+4+2*4), der so aussehen könnte:

\textbf{Beispiel für einen Datenstrom}

VN\textasciitilde \textasciitilde \textasciitilde NN\textasciitilde\textasciitilde12345678

Datenströme sind fast immer 1920 Bytes lang.


\subsection{Projektziel}
\label{sec:Projektziel}
Es soll ein Programm geschrieben werden, das nach Eingabe von Datenstrom und 1920Schema den Datenstrom dem Schema entsprechend zergliedert und anzeigt. Das Programm wird im Folgenden 1920Parser genannt.

\textbf{Gewünschte Ausgabe} (mit den Beispielwerten der Abschnitte 1.2.1 und 1.2.2)

\begin{tabularx}{0px}{l}
01 Daten\\
\ \ \ 03 Personendaten\\
\ \ \ \ \ \ 05 Vorname=VN\textasciitilde \textasciitilde \textasciitilde\\
\ \ \ \ \ \ 05 Nachname=NN\textasciitilde \textasciitilde\\
\ R03 Gesamter-Name=VN\textasciitilde \textasciitilde \textasciitilde NN\textasciitilde \textasciitilde\\
\ \ \ 03 Bestellungen(1)\\
\ \ \ \ \ \ 05 Artikelnr=1234\\
\ \ \ 03 Bestellungen(2)\\
\ \ \ \ \ \ 05 Artikelnr=5678\\
\end{tabularx}

\subsection{Projektbegründung}
\label{sec:Projektbegruendung}
Bei Kundenreklamationen, Änderungen an Programmen und Neuentwicklungen stehen die Warehouse-Programmierer oft vor den Problemen:

\begin{itemize}
\item Wert einer Schemadatei-Variablen in einem Datenstrom finden.
\item Datenstrom-Bytes einer Schemadatei-Variablen zuordnen.
\end{itemize}

Gegenwärtig zählen die Programmierer die passende Anzahl von Bytes in Schema und Datenstrom ab, einige erfahrene kennen die wichtigsten Schemadateien auch teilweise auswendig. 


\subsection{Projektschnittstellen}
\label{sec:Projektschnittstellen}
Benutzer des Programms sind die Programmierer der Phoenix Group IT GmbH, hauptsächlich die aus der Abteilung Warehouse.

1920Parser soll nicht unmittelbar mit anderen Systemen interagieren. Vorgesehen ist, dass die Benutzer die notwendigen Angaben in eine Eingabemaske hineinkopieren.

Projektgenehmigung und die Bereitstellung von Ressourcen erfolgt durch die Ausbildende Frau Birgit Günther, die Projektbetreuung und die Abnahme des Programms durch Herrn Marco Kemmer. Herr Kemmer arbeitet in der Abteilung Warehouse als COBOL-Entwickler, er kommuniziert die Kundenwünsche und will das Programm in Zukunft auch selbst nutzen.





% !TEX root = ../Projektdokumentation.tex
\section{Projektplanung}
\label{sec:Projektplanung}


\subsection{Projektphasen}
\label{sec:Projektphasen}
Für das Projekt standen 70 Stunden zur Verfügung. Es findet im Zeitraum vom 11.04.2016 - 15.05.2016 statt.

\paragraph{Beispiel}
Tabelle~\ref{tab:Zeitplanung} zeigt die grobe Zeitplanung für das Projekt.
\tabelle{Zeitplanung}{tab:Zeitplanung}{ZeitplanungKurz}\\
Eine detailliertere Zeitplanung findet sich im \Anhang{app:Zeitplanung}.

\subsection{Ressourcenplanung}
\label{sec:Ressourcenplanung}
Zur Projektdurchführung werden folgende Ressourcen eingesetzt:
Windows 7, Visual Studio Professional 2010, Microsoft Visio 2010, TexMaker, XMLSpy 2007, OpenText HostExplorer, Büro mit PC mit Verbindung zum Mainframe und Verbindung zum  Internet, Strom, Projektbetreuer


\subsection{Entwicklungsprozess}
\label{sec:Entwicklungsprozess}
Es wurde ein agiler Entwicklungsprozess angewendet, der an Extreme Programming angelehnt war. Der im Projekt angewandte Prozess beinhaltete die Extreme Programming-Praktiken, testgetriebene Entwicklung, häufige Kundeneinbeziehung, häufiges Refaktorisieren, kurze Iterationen und einfaches Design. 







% !TEX root = ../Projektdokumentation.tex
\section{Analysephase} 
\label{sec:Analysephase}

\subsection{Ist-Analyse} 
\label{sec:IstAnalyse}
Die Ist-Analyse fand durch Befragung von Herrn Kemmer statt. 
Die Entwickler der Abteilung Warehouse müssen bei Kundenreklamationen, Programmänderungen und Neuentwicklungen die Werte in den beschriebenen Datenströme analysieren. Diese stammen aus den Logdateien des zentralen Lagerrechners, aus der Eingabemaske des Mainframes oder auch aus erhaltenen Email-Anhängen. Um die Bedeutung der Werte im Datenstrom herauszufinden, zählen die Entwickler gegenwärtig die Bytes in Schema und Datenstrom ab, einige erfahrenere kennen die wichtigsten Schemadateien auch zum Teil auswendig.
Im Warehouse arbeiten etwa 20 Entwickler, etwa 10 davon haben regelmäßig mit solchen Datenströmen zu tun. Es werden Dutzende verschiedene 1920Schemas für verschiedenste Zwecke verwendet, aber mit nur etwa 10 arbeiten die Entwickler regelmäßig.


\subsection{Wirtschaftlichkeitsanalyse}
\label{sec:Wirtschaftlichkeitsanalyse}
Das Programm soll den Entwicklern das bisher notwendige, fehleranfällige Abzählen von Zeichen in Schema und Datenstrom abnehmen. Das Projekt verspricht dadurch nicht nur, den Entwicklern Zeit zu sparen, sondern auch Zählfehler wirkungsvoll zu verhindern.

\subsubsection{\gqq{Make or Buy}-Entscheidung}
\label{sec:MakeOrBuyEntscheidung}
Die Anforderungen sind so speziell, dass fast auszuschließen ist, dass es außerhalb von Phoenix ein Programm gibt, das die Anforderungen erfüllt. Zu bemerken ist aber, dass viele Programme bei Phoenix mit Schemadateien arbeiten. Der Autor fragte deshalb nach, ob Phoenix eventuell schon ein Programm hat, das die Anforderungen erfüllt. Der Auftraggeber antwortete, dass die Datenströme aus so unterschiedlichen Umgebungen stammen (Unix-Lagerrechner, Mainframe, Email), dass die COBOL-Programme, die sie verarbeiten, nur eingeschränkt eingesetzt werden können. Es wurde entschieden, das Programm selbst neu zu schreiben.

\subsubsection{Projektkosten}
\label{sec:Projektkosten}
Die Kosten für die Durchführung des Projekts setzen sich sowohl aus Personal-, als auch aus Ressourcenkosten zusammen.
Der Projektersteller ist Umschüler und erhält deshalb von seinem Ausbildungsbetrieb keine Vergütung.

\begin{eqnarray}
7,7 \mbox{ h/Tag} \cdot 220 \mbox{ Tage/Jahr} = 1694 \mbox{ h/Jahr}\\
\eur{0}\mbox{/Monat} \cdot 13,3 \mbox{ Monate/Jahr} = \eur{0} \mbox{/Jahr}\\
\frac{\eur{0} \mbox{/Jahr}}{1694 \mbox{ h/Jahr}} = \eur{0,00}\mbox{/h}
\end{eqnarray}

Dadurch ergibt sich also ein Stundenlohn von \eur{0,00}
Die Durchführungszeit des Projekts beträgt 70 Stunden. Für die Nutzung von Ressourcen\footnote{Räumlichkeiten, Arbeitsplatzrechner etc.} wird 
ein pauschaler Stundensatz von \eur{12} angenommen. Für die anderen Mitarbeiter wird pauschal ein Stundenlohn von \eur{23} angenommen. 
Eine Aufstellung der Kosten befindet sich in Tabelle~\ref{tab:Kostenaufstellung} und sie betragen insgesamt \eur{1015,00}.
\tabelle{Kostenaufstellung}{tab:Kostenaufstellung}{Kostenaufstellung.tex}


\subsubsection{Amortisationsdauer}
\label{sec:Amortisationsdauer}
Es wird davon ausgegangen, dass ausschließlich Entwickler der Abteilung Lager das Programm nutzen werden. Nach Einschätzung von Herrn Kemmer arbeitet die Hälfte der 20 Entwickler regelmäßig mit Schemadateien und dass das Programm jedem täglich 10 Minuten einsparen kann.
Bei einer Einsparung von 10 Minuten am Tag für 10 Entwickler an 220 Arbeitstagen im Jahr ergibt sich eine gesamte Zeiteinsparung von 
\begin{eqnarray}
10 \cdot 220 \mbox{ Tage/Jahr} \cdot 10 \mbox{ min/Tag} = 22000 \mbox{ min/Jahr} \approx 366,67 \mbox{ h/Jahr} 
\end{eqnarray}

Dadurch ergibt sich eine jährliche Einsparung von 
\begin{eqnarray}
366,67 \mbox{h} \cdot \eur{(23 + 12)}{\mbox{/h}} = \eur{12833,45}
\end{eqnarray}

Die Amortisationsdauer beträgt also $\frac{\eur{1015,00}}{\eur{12833,45}\mbox{/Jahr}} \approx 0,08 \mbox{ Jahre} \approx 1 \mbox{ Monat}$.


\subsection{Anwendungsfälle}
\label{sec:Anwendungsfaelle}
Der mit Abstand wichtigste Anwendungsfall ist, dass das Programm nach Angabe von Datenstrom und Schema den Datenstrom entsprechend dem Schema zergliedert anzeigt.
Der Kunde nannte noch einige Wünsche mehr zur Funktionalität. Abbildung \ref{app:Usecase} stellt die Anwendungsfälle in einem Usecase-Diagramm dar, sie befindet sich im Anhang.


\subsection{Qualitätsanforderungen}
\label{sec:Qualitaetsanforderungen}
Schemas müssen frei angebbar sein und Datenströme richtig zergliedert werden.
Beim Abnahmetest soll dies anhand der 5 wichtigsten Schemadateien überprüft werden.
Die Performance des Programmes ist ziemlich egal, es soll aber flüssig benutzbar sein (Zergliederung von Datenstrom und Anzeige in unter 5 Sekunden).
Weil die Benutzer Profis sind, ist es nicht unbedingt notwendig, für alles eine Eingabemaske bereitzustellen. Es genügt auch, wenn eine Einstellung durch Editieren einer Konfigurations-Datei geändert werden kann.


\Zwischenstand{Analysephase}{Analyse}


% !TEX root = ../Projektdokumentation.tex
\section{Entwurfsphase} 
\label{sec:Entwurfsphase}

\subsection{Zielplattform}
\label{sec:Zielplattform}
Das Programm soll auf den Entwicklerrechnern der Phoenix laufen. Intel i5 Prozessoren mit 4 GB Arbeitsspeicher.
Auf den Entwicklerrechnern laufen 32 Bit Version von Windows 7. Phoenix programmiert in COBOL, C++ und C\#. COBOL fiel für ein Windows-Tool aus. Die Wahl fiel auf C\# (Garbage Collection, moderner GUI-Designer).


\subsection{Architekturdesign}
\label{sec:Architekturdesign}
GUI und Anwendungslogik wurden getrennt.


\subsection{Entwurf der Benutzeroberfläche}
\label{sec:Benutzeroberflaeche} 
\begin{itemize}
	\item Entscheidung für die gewählte Benutzeroberfläche (\zB GUI, Webinterface).
	\item Beschreibung des visuellen Entwurfs der konkreten Oberfläche (\zB Mockups, Menüführung).
	\item \Ggfs Erläuterung von angewendeten Richtlinien zur Usability und Verweis auf Corporate Design.
\end{itemize}

\paragraph{Beispiel}
Beispielentwürfe finden sich im \Anhang{app:Entwuerfe}.



\subsection{Geschäftslogik}
\label{sec:Geschaeftslogik}
Zunächst wird das Schema in eine Baumstruktur überführt.



\paragraph{Beispiel}
Ein Klassendiagramm, welches die Klassen der Anwendung und deren Beziehungen untereinander darstellt kann im \Anhang{app:Klassendiagramm} eingesehen werden.

\Abbildung{Modulimport} zeigt den grundsätzlichen Programmablauf beim Einlesen eines Moduls als \ac{EPK}.
\begin{figure}[htb]
\centering
\includegraphicsKeepAspectRatio{modulimport.pdf}{0.9}
\caption{Prozess des Einlesens eines Moduls}
\label{fig:Modulimport}
\end{figure}


\subsection{Maßnahmen zur Qualitätssicherung}
\label{sec:Qualitaetssicherung}
\begin{itemize}
	\item Welche Maßnahmen werden ergriffen, um die Qualität des Projektergebnisses (siehe Kapitel~\ref{sec:Qualitaetsanforderungen}: \nameref{sec:Qualitaetsanforderungen}) zu sichern (\zB automatische Tests, Anwendertests)?
	\item \Ggfs Definition von Testfällen und deren Durchführung (durch Programme/Benutzer).
\end{itemize}


\subsection{Pflichtenheft/Datenverarbeitungskonzept}
\label{sec:Pflichtenheft}
\begin{itemize}
	\item Auszüge aus dem Pflichtenheft/Datenverarbeitungskonzept, wenn es im Rahmen des Projekts erstellt wurde.
\end{itemize}

\paragraph{Beispiel}
Ein Beispiel für das auf dem Lastenheft (siehe Kapitel~\ref{sec:Lastenheft}: \nameref{sec:Lastenheft}) aufbauende Pflichtenheft ist im \Anhang{app:Pflichtenheft} zu finden.


\Zwischenstand{Entwurfsphase}{Entwurf}

% !TEX root = ../Projektdokumentation.tex
\section{Implementierungsphase} 
\label{sec:Implementierungsphase}
Zuerst legte der Autor ein neues SVN-Repository für sein Programm an.
Die Implementierung begann mit dem Anlegen eines neuen Repositories mit der vorgegebenen Verzeichnisstruktur (branch, tag, trunk) in SVN und dem Erstellen einer neuen Solution in Microsoft Visual Studio 2010. Der Solution wurde ein NUnit-Testprojekt für die Unit-Tests hinzugefügt.

\subsection{Implementierung der Datenstrukturen}
\label{sec:ImplementierungDatenstrukturen}
Die hierarchische Gliederung der Schemavariablen wird im Programm durch eine Baumstruktur abgebildet. Schemavariablen können Gruppen- oder Wertevariablen sein, Entsprechend gibt es im Programm die Klassen GroupNode und ValueNode. Die  beiden Klassen haben gemeinsamen Attribute und Methoden, unter anderem
\begin{itemize}
\item Level
\item VarName
\item RepeatCount
\item RepeatIndex
\item Comment
\item AssignValue()
\item ToString()
\end{itemize}. Diese wurden in der Klasse AbstractNode definiert und von Group- und ValueNode geerbt.
GroupNodes erhielten ein Attribut children vom Typ List<AbstractNode> und können so wiederum auf weitere GroupNodes verweisen (und diese auch wieder). Dadurch kann eine beliebig tief verschachtelbare Baumstruktur repräsentiert werden. 


Die Methode, die das Schema parst, soll den Wurzelknoten als GroupNode zurückgeben. Daher musste die Implementierung mit den Klassen der Baumstruktur (AbstractNode, GroupNode und ValueNode) begonnen werden. Zum Schreiben von Unit-Tests wurde eine Methode zum vergleichen der Ergebnisse benötigt. Dazu wurden die Methoden Equals() und getHashCode() der Klasse object überschrieben. Danach wurde die Baumstruktur gemockt und die Methoden AssignValue() und ToString() implementiert.

Danach wurde die Methode Parse() der Klasse Schema geschrieben. Hier war ein Problem, dass gleichzeitig 


\subsection{Implementierung der Benutzeroberfläche}
\label{sec:ImplementierungBenutzeroberflaeche}
Die \ac{GUI} muss dem Benutzer Funktionalität bereitstellen, damit er Schema und Datenstrom angeben kann und sie muss den zergliederten Datenstrom anzeigen können.
Um diese Funktionalitäten zu bieten werden 3 Textboxen angezeigt, jeweils mit einem Label als Überschrift. Auf ein Menü wurde zunächst verzichtet.

Die Angaben in den Schemadateien sind formatiert, gleiche Angaben wie Level oder Variablenname stehen untereinander. Damit die Angaben in der GUI auch untereinander stehen, wurde für das Schema-Textfeld ein Monospace-Font verwendet (Courier New). Um das Erscheinungsbild einheitlich zu halten, zeigen auch die Textfelder für Datenstrom und Ergebnis Ihren Inhalt mit diesem Zeichensatz an.
\paragraph{Beispiel}
Screenshots der Anwendung in der Entwicklungsphase mit Dummy-Daten befinden sich im \Anhang{Screenshots}.


\subsection{Implementierung der Geschäftslogik}
\label{sec:ImplementierungGeschaeftslogik}

\subsubsection{Grundschema der rekursive Methoden von AbstractNode/GroupNode}
Group- und ValueNode erben von AbstractNode Methoden. GroupNode definiert definiert diese Methoden rekursiv. Alle diese Methoden funktionieren nach dem gleichen Prinzip:

GroupNode ruft für jedes seiner Kinder die gleiche Methode erneut auf (Rekursion). Aus den Rückgabewerten der Kinder wird ein Wert akkumuliert und dieser zurückgegeben.
Eine ValueNode hat keine Kinder und kann direkt ein Wert zurückgeben. Bei ValueNodes und GroupNodes ohne Kindern enden die Rekursionen.



\subsubsection{Parsen des Schemas in eine Baumstruktur}
Da der Baum einen einzigen Wurzelknoten\footnote{Knoten, der Vorfahr von allen anderen Knoten ist} haben soll, Schemas aber nicht zwingend einen haben, erhält der Baum einen künstlichen Wurzelknoten mit der Stufennummer 0. Würde man keinen solchen Knoten einführen, müsste statt mit GroupNode mit List<AbstractNode> als Rückgabewert gearbeitet werden, da es auf höchster Ebene mehrere Geschwister geben kann oder Wertvariablen ohne Eltern.

Eine Schwierigkeit des Algorithmus war, dass zu jedem Zeitpunkt in mehrere Richtungen gegangen werden muss. Zum einen muss die nächste Zeile des Schemas in einen Knoten verwandelt werden. Ein Knoten kann sich aber wiederholen.
Hier hat der Autor einige Zeit verbraucht mit rekursiven Ideen.

Der Algorithmus zum Parsen des Schemas nutzt einen Stack. Der Stack enthält zu jedem Zeitpunkt die Vorfahren des gerade bearbeiteten Knotens, mit seinem Elternknoten oben und dem Wurzelknoten unten. x wird mit seinem





\subsubsection{Zuweisen der Werte aus dem Datenstrom}
Die Methode zum Zuweisen eines Wertes an einen Knoten hat den Prototypen 

int AssignValue(string data).

Die Methode erwartet den Datenstrom als string und gibt die Anzahl der verbrauchten Bytes zurück. Das aufrufende Objekt (normalerweise eine GroupNode) bekommt dadurch Informationen, wie viele Bytes des Datenstroms benutzt wurden. So kann der Anfang des Datenstroms für jedes Kind passend verschoben werden (und bei Redefines auch wieder zurückverschoben werden).

ValueNodes und GroupNodes machen beim Aufruf Folgendes:

ValueNodes weisen ihrem Value-Attribut die benötigte Anzahl Buchstaben aus data zu und geben diese Anzahl zurück.


GroupNodes rufen für alle ihre Kinder nacheinander AssignValue(data) auf. Jedes Kind gibt die Anzahl der verwendeten Bytes zurück. GroupNodes merken sich diese Anzahl und können dadurch für jedes Kind den Anfang des Datenstroms um diese Anzahl verschieben. GroupNode gibt die Summe der verwendeten Bytes aller seiner Kinder zurück.

\paragraph{Beispiel}
Die Klasse \texttt{Com\-par\-ed\-Na\-tu\-ral\-Mo\-dule\-In\-for\-ma\-tion} findet sich im \Anhang{app:CNMI}.  


\Zwischenstand{Implementierungsphase}{Implementierung}





% !TEX root = ../Projektdokumentation.tex
\section{Abnahmephase} 
\label{sec:Abnahmephase}
Der Abnahmetest erfolgte durch den Projektbetreuer Herrn Kemmer. Herr Kemmer ließ sich die erstellten Unit-Tests zeigen überprüfte, dass alle Tests erfolgreich waren. Danach ludt er die 5 am Häufigsten verwendeten Schemas zusammen mit passenden Datenströmen vom Mainframe. 1920Parser zeigte alle getesteten Schemas richtig an und zergliederte die Datenströme in gewünschter Weise. Im Anschluss testete Herr Kemmer die Funktionalität zum Speichern und zur automatischen Auswahl von Schemas. Nach Löschung der XML-Konfigurationsdatei trat bei der Auswahl des Menüpunktes "Config editieren" eine FileNotFoundException auf.

\Zwischenstand{Abnahmephase}{Abnahme}

% !TEX root = ../Projektdokumentation.tex
\section{Einführungsphase}
\label{sec:Einfuehrungsphase}
Die Einführung von 1920Parser wurde von Herrn Kemmer übernommen. Das Programm soll im Intranet bereitgestellt werden, da die Benutzer nicht alle \ac{SVN} installiert haben.

%\Zwischenstand{Einführungsphase}{Einfuehrung}

% !TEX root = ../Projektdokumentation.tex
\section{Dokumentation}
\label{sec:Dokumentation}
Für 1920Parser wurde eine Benutzerdokumentation und eine Entwicklerdokumentation geschrieben.

\paragraph{Beispiel}
Ein Ausschnitt aus der erstellten Benutzerdokumentation befindet sich im \Anhang{app:BenutzerDoku}.
Die Entwicklerdokumentation wurde mittels PHPDoc\footnote{Vgl. \cite{phpDoc}} automatisch generiert. Ein beispielhafter Auszug aus der Dokumentation einer Klasse findet sich im \Anhang{app:Doc}. 

\Zwischenstand{Dokumentation}{Dokumentation}

% !TEX root = ../Projektdokumentation.tex
\section{Fazit} 
\label{sec:Fazit}

\subsection{Soll-/Ist-Vergleich}
\label{sec:SollIstVergleich}
Wie gewünscht können im Programm Datenstrom und Schema frei angegeben werden. Das Ergebnis wird richtig angezeigt.
Insbesondere das Zuweisen der Werte aus dem Datenstrom und die Ausgabe des Baumes als String nahmen deutlich weniger Zeit in Anspruch, als ursprünglich geplant. So konnte noch der Anwendungsfall ``Schema speichern'' im Rahmen des Projektes umgesetzt werden.
Herr Kemmer ist mit dem Programm sehr zufrieden.
\tabelle{Soll-/Ist-Vergleich}{tab:Vergleich}{Zeitnachher.tex}
\subsection{Lessons Learned}
\label{sec:LessonsLearned}
RichTextBox ist besser als TextBox.
Interessant war, wie sich der Projektumfang erweitert hatte (ursprünglich war der Vorschlag, dass nur ein bestimmtes Schema zergliedert werden soll)
Rekursion ist nicht gut darstellbar mit UML. Für Polymorphie gilt das selbe. Stacks sind super. Rekursion vereinfacht manche Aufgaben enorm. Bemerkenswert, wie sehr sich die Anforderungen ausgeweitet haben.
Bemerkenswert war, dass sich der Projektumfang stark erhöhte (ursprünglich wurde mir der Projektvorschlag gemacht, nur die Schemadatei VK60 zu zergliedern).
\subsection{Ausblick}
\label{sec:Ausblick}
Im Rahmen des Projektes wurden nicht alle Wunschkriterien erfüllt. Der Autor hofft aber, dass diese Aufgabe vielleicht an nachfolgende Auszubildende übergeben wird und sich 1920Parser noch weiterentwickelt. Vielleicht wird einmal die Größe der Datenströme geändert, dann könnte der Algorithmus von 1920Parser verwendet werden.







% Literatur ------------------------------------------------------------------
\clearpage
\renewcommand{\refname}{Literaturverzeichnis}
\bibliography{Bibliographie}
\bibliographystyle{Allgemein/natdin} % DIN-Stil des Literaturverzeichnisses
% !TEX root = Projektdokumentation.tex
\clearpage
\addsec{Eidesstattliche Erklärung}

% Hinweis: die eidesstattliche Erklärung ist ggfs. an die Richtlinie der IHK anzupassen

Ich, \autorName, versichere hiermit, dass ich meine \textbf{\betreff} mit dem
Thema
\begin{quote}
\textit{\kompletterTitel}
\end{quote}
selbständig verfasst und keine anderen als die angegebenen Quellen und Hilfsmittel benutzt habe,
wobei ich alle wörtlichen und sinngemäßen Zitate als solche gekennzeichnet habe. Die Arbeit
wurde bisher keiner anderen Prüfungsbehörde vorgelegt und auch nicht veröffentlicht.\\[6ex]

\abgabeOrt, den \abgabeTermin


\rule[-0.2cm]{5.5cm}{0.5pt}

\textsc{\autorName}


% Anhang ---------------------------------------------------------------------
\def \renePage {iiix}
\clearpage
\appendix
\pagenumbering{roman}
% !TEX root = Projektdokumentation.tex
\section{Anhang}
\subsection{Detaillierte Zeitplanung}
\label{app:Zeitplanung}

\tabelleAnhang{ZeitplanungKomplett}

\clearpage
\subsection{Beispiel für ein echtes 1920Schema}
\label{app:RealSchema}
\begin{figure}[htb]
\centering
\includegraphicsKeepAspectRatio{schemabim.pdf}{0.74}
\caption{Beispiel für ein echtes 1920Schema}
\end{figure}
\clearpage

\subsection{Use Case-Diagramm}
\label{app:UseCase}
\begin{figure}[htb]
\centering
\includegraphicsKeepAspectRatio{reneusecase.pdf}{1.0}
\caption{Use Case-Diagramm}
\end{figure}
\clearpage

\subsection{Pflichtenheft (Auszug)}
\label{app:Pflichtenheft}
Es folgt ein Auszug aus dem Lastenheft mit Fokus auf die Anforderungen:

Die Anwendung muss folgende Anforderungen erfüllen: 
\begin{enumerate}[itemsep=0em,partopsep=0em,parsep=0em,topsep=0em]
\item Muss-Kriterien:
	\begin{enumerate}
	\item Datenstrom und Schema müssen frei angegeben werden können.
	\item Das Programm muss den Datenstrom anhand des Schemas zergliedert anzeigen.
	\end{enumerate}
\item Soll-Kriterien:
	\begin{enumerate}
	\item Benutzer sollen ein eingegebenes Schemas speichern können.
	\item Benutzer sollen gespeicherte Schemas laden können.
	\end{enumerate}
\item Kann-Kriterien:
	\begin{enumerate}
    \item Anhand des Transaktionscodes im Datenstrom soll das richtige Schema - falls vorhanden - automatisch ausgewählt werden.
    \item Der Benutzer soll ein Kriterium angeben können, anhand dessen die passenden Datenströme aus der Logdatei ausgewählt werden.
    \item Wenn das Länge-Feld eine Angabe zur Anzahl der Nachkommastellen hat, soll das Programm in der Ergebnis-Anzeige an der richtigen Stelle ein Komma einfügen.

    \item Variablenzeilen sollen ausgeblendet werden können.
        \end{enumerate}
\item Ausschluss-Kriterien:
	\begin{enumerate}
    \item Das Programm soll Schema und Datenstrom nicht auf Korrektheit oder Plausibilität prüfen
(z. B. doppelte Variablennamen, Redefines ohne Vorgänger, ungültige Stufennummern, Bytelänge 0, Char mit Nachkommastellen).
    \end{enumerate}
\end{enumerate}

\clearpage

\subsection{Klassendiagramm}
\label{app:classd}
\begin{figure}[htb]
\centering
\includegraphicsKeepAspectRatio{classd.pdf}{0.74}
\caption{Klasseniagramm}
\end{figure}
\clearpage

\subsection{Aktivitätsdiagramm - Parsen eines 1920Schemas}
\label{app:AktivitaetsdiagrammSchemaParsen}
\begin{figure}[htb]
\centering
\includegraphicsKeepAspectRatio{aktivitaet_createTree11.pdf}{0.9}
\caption{Aktivitätsdiagramm - Parsen eines 1920Schemas}
\end{figure}
\clearpage


\subsection{Objektdiagramm - Parsen des Schemas}
\label{app:ObjDiagramsdiagrammSchemaParsen}
\begin{figure}[htb]
\centering
\includegraphicsKeepAspectRatio{objectdiagram.pdf}{0.74}
\caption{Objektdiagramm nach dem Parsen und dem Zuweisen der Werte}
\end{figure}
\clearpage




\subsection{Quelltext Schema.cs}
\label{app:SchemaSrc}
\begin{figure}[htb]
\centering
\includegraphicsKeepAspectRatio{schemasrc.pdf}{0.74}
\caption{Schema.cs}
\end{figure}
\clearpage




\subsection{Quelltext GroupNode.cs}
\label{app:GroupNodeSrc}
\begin{figure}[htb]
\centering
\includegraphicsKeepAspectRatio{groupnodesrc.pdf}{0.74}
\caption{GroupNode.cs}
\end{figure}
\clearpage


\subsection{Quelltext ValueNode.cs}
\label{app:ValueNodeSrc}
\begin{figure}[htb]
\centering
\includegraphicsKeepAspectRatio{valuenodesrc.pdf}{0.74}
\caption{ValueNode.cs}
\end{figure}
\clearpage


\subsection{Entwicklerdokumentation}
\label{app:Entwicklerdokumentation}
\begin{figure}[htb]
\centering
\includegraphicsKeepAspectRatio{developerdocumentation.pdf}{1}
\caption{Entwicklerdokumentation}
\end{figure}
\clearpage


\end{document}
